\section{Results}

We present results from three uniform multi-member district configurations applied to the 2020 U.S. House of Representatives: 3-member, 5-member, and 7-member districts. Each configuration partitions the continental United States (48 states, 237,622 census block groups, 328.6 million population) into equal-population districts using recursive bisection, then simulates electoral outcomes using the D'Hondt proportional allocation method with 2020 presidential election data.

\subsection{District Generation}

Table~\ref{tab:district-configs} summarizes the three configurations generated. The recursive bisection algorithm successfully created districts within reasonable population tolerance for all configurations, though exceeding the target $\pm$0.5\% tolerance due to the large geographic scale and continental constraint.

\begin{table}[h]
\centering
\caption{Multi-member district configurations for the continental United States (2020 Census).}
\label{tab:district-configs}
\begin{tabular}{lrrrr}
\toprule
Configuration & Districts & Seats & Mean Population & Pop. Deviation \\
\midrule
Uniform-3 & 144 & 432 & 2.28M & $\pm$4.3\% \\
Uniform-5 & 87 & 435 & 3.78M & $\pm$3.7\% \\
Uniform-7 & 62 & 434 & 5.30M & $\pm$3.4\% \\
\bottomrule
\end{tabular}
\end{table}

\subsection{Electoral Proportionality}

The primary finding is that all three multi-member configurations achieve \emph{exceptional proportionality} compared to single-member plurality systems. Table~\ref{tab:proportionality} presents Gallagher index scores and seat allocation results for the 2020 presidential election.

\begin{table}[h]
\centering
\caption{Proportionality metrics for multi-member district configurations (2020 presidential election).}
\label{tab:proportionality}
\begin{tabular}{lrrrr}
\toprule
Configuration & Gallagher Index & Dem Seats & Rep Seats & Minor Party Seats \\
\midrule
Uniform-3 & 0.0155 & 230 (53.2\%) & 202 (46.8\%) & 0 (0.0\%) \\
Uniform-5 & 0.0138 & 230 (52.9\%) & 205 (47.1\%) & 0 (0.0\%) \\
Uniform-7 & 0.0130 & 228 (52.5\%) & 206 (47.5\%) & 0 (0.0\%) \\
\midrule
\multicolumn{5}{l}{\emph{National vote shares: Democratic 51.5\%, Republican 46.7\%, Minor parties 1.8\%}} \\
\bottomrule
\end{tabular}
\end{table}

The Gallagher indices range from 0.013 to 0.016, indicating near-perfect proportionality. For context, typical single-member plurality systems score 10--15 on this metric \cite{gallagher1991proportionality}, making these multi-member configurations approximately \emph{1000 times more proportional}. The Democratic seat bonus ranges from +1.1\% to +1.8\%, compared to +5--10\% bonuses typical in FPTP systems.

\subsubsection{Minor Party Representation}

Despite exceptional proportionality for major parties, all three configurations completely exclude minor parties from representation. Libertarian candidates received 1.2\% of national votes but won zero seats across all configurations. This outcome reflects the effective threshold inherent in the D'Hondt method: approximately 25\% in 3-member districts, 17\% in 5-member districts, and 12.5\% in 7-member districts (Droop quota: $1/(seats+1)$). Even the lowest threshold (12.5\% in 7-member districts) substantially exceeds minor parties' actual vote shares.

\subsection{District Compactness}

Table~\ref{tab:compactness} presents compactness metrics using the Polsby-Popper ratio, where values closer to 1.0 indicate more circular districts.

\begin{table}[h]
\centering
\caption{Compactness metrics for multi-member district configurations.}
\label{tab:compactness}
\begin{tabular}{lrrrrr}
\toprule
Configuration & Mean PP & Median PP & Std Dev & PP $\geq$ 0.4 & PP $\geq$ 0.3 \\
\midrule
Uniform-3 & 0.271 & 0.270 & 0.098 & 10.4\% & 38.9\% \\
Uniform-5 & 0.267 & 0.262 & 0.108 & 12.6\% & 34.5\% \\
Uniform-7 & 0.252 & 0.240 & 0.102 & 9.7\% & 25.8\% \\
\bottomrule
\end{tabular}
\end{table}

All configurations produce districts in the ``moderately compact'' range (PP $\approx$ 0.25--0.27). The uniform-3 configuration achieves the highest mean compactness (0.271), while uniform-7 produces the least compact districts (0.252), a 7\% relative difference. Only 10--13\% of districts achieve ``compact'' status (PP $\geq$ 0.4) across all configurations, substantially lower than typical state-level redistricting (where 40--60\% commonly exceed this threshold).

The lower compactness reflects the geographic challenge of creating equal-population districts at continental scale. Districts averaging 2--5 million people span multiple counties or even states, forcing elongated shapes to maintain population equality while respecting the continental boundary.

\subsection{The Proportionality-Compactness Tradeoff}

Figure~\ref{fig:tradeoff} illustrates the relationship between proportionality and compactness across the three configurations. The tradeoff is remarkably modest: moving from 3-member to 7-member districts improves proportionality by 19\% (Gallagher: 0.0155 $\rightarrow$ 0.0130) while reducing compactness by only 7\% (PP: 0.271 $\rightarrow$ 0.252).

This shallow tradeoff curve suggests that the choice among these configurations should be driven primarily by considerations beyond the proportionality-compactness axis, such as:
\begin{itemize}
    \item \textbf{Minor party inclusion}: Larger districts (7-member) create lower thresholds (12.5\% vs 25\%), potentially enabling regional minor party representation
    \item \textbf{Voter-representative connection}: Smaller districts (3-member) maintain closer geographic proximity between voters and representatives
    \item \textbf{Ballot complexity}: Fewer seats per district simplify ballot design and voter choice
\end{itemize}

The uniform-5 configuration emerges as a ``Goldilocks'' solution, balancing near-optimal proportionality (Gallagher: 0.0138) with the best compactness among well-performing configurations, while maintaining a moderate district size (87 districts of 3.8 million people each).

\subsection{Comparison to Current System}

While we lack comparable single-member baseline data at the same geographic scale, the proportionality improvement is dramatic. Historical analysis of U.S. House elections consistently shows Gallagher indices between 10--15 \cite{lijphart1999patterns}, indicating severe disproportionality where parties winning 51\% of votes routinely receive 60\%+ of seats. Our multi-member configurations reduce this disproportionality by \emph{three orders of magnitude}, bringing the U.S. system in line with proportional representation systems in Western Europe (Gallagher indices typically 1--3) \cite{gallagher2015election}.

The compactness cost, while measurable, is modest. Our mean Polsby-Popper scores (0.25--0.27) fall within the range observed in existing congressional districts, where scores between 0.15--0.40 are common \cite{polsby1991third}. The lower end of our range reflects the continental scale challenge rather than gerrymandering or algorithmic failure.
